%%%%%%%%%%%%%%%%%%%%%%%%%%%%%%%%%%%%%%%%%%%%%%%%%%%%%%%%%%%%%%%%%%%%%%
% Lớp 10 | Véctơ (chưa xét tọa độ) | Tổng và hiệu của hai véctơ
% Dạng 0H5?2-2: Tính đúng-sai của 1 đẳng thức véctơ
% Số câu: 23  (sinh tự động bởi tach_cau.py; từ nay thầy sửa trực tiếp file này)
%%%%%%%%%%%%%%%%%%%%%%%%%%%%%%%%%%%%%%%%%%%%%%%%%%%%%%%%%%%%%%%%%%%%%%

%%%%%%%%%%%%%%%% TRẮC NGHIỆM 4 PHƯƠNG ÁN %%%%%%%%%%%%%%%%

% nguồn: MaHoaNH_TracNghiem_Lop10
\begin{ex}%[0H5N2-2]%[ma:6eda98d4]
	Cho tam giác $ABC$. Khẳng định nào sau  đây là \textbf{sai}?
	 \choice
		 {$\overrightarrow{AB}-\overrightarrow{AC}=\overrightarrow{CB}$}
		 {$\overrightarrow{AB}+\overrightarrow{BA}=\overrightarrow{0}$}
		{\True $\overrightarrow{AB}-\overrightarrow{AC}=\overrightarrow{BC}$} 
		 {$\overrightarrow{AB}+\overrightarrow{BC}=\overrightarrow{AC}$}
	 \loigiai{
	Ta có $\overrightarrow{AB}-\overrightarrow{AC}=\overrightarrow{CB}$ nên khẳng định $\overrightarrow{AB}-\overrightarrow{AC}=\overrightarrow{BC}$ là sai.} 
\end{ex}

% nguồn: MaHoaNH_TracNghiem_Lop10
\begin{ex}%[0H5N2-2]%[ma:6f06b4df]
	Cho $M$ là trung  điểm của đoạn thẳng $AB$. Khẳng định nào sau đây là khẳng định đúng?
	\choice 
		{$\overrightarrow{IA}+\overrightarrow{IB}=\overrightarrow{AB}$  với $I$ là điểm bất kì}
		{\True $\overrightarrow{AM}+\overrightarrow{BM}=\overrightarrow{0}$} 
		{$\overrightarrow{IA}+\overrightarrow{IB}=\overrightarrow{IM}$ với $I$ là điểm bất kì} 
		 {$\overrightarrow{AM}+\overrightarrow{MB}=\overrightarrow{0}$}
	 \loigiai{
	Vì $M$ là trung  điểm của $AB$ nên ta luôn có $\overrightarrow{AM}+\overrightarrow{BM}=\overrightarrow{0}$.}
\end{ex}

% nguồn: MaHoaNH_TracNghiem_Lop10
\begin{ex}%[0H5N2-2]%[ma:3d1391b2]
	Cho ba điểm phân biệt $A$, $B$, $C$.  Mệnh đề nào sau đây \textbf{đúng}?
	 \choice
		 {$AB+BC=AC$}
		{\True $\vv{AB}+\vv{BC}+\vv{CA}=\vv{0}$} 
		 {$\vv{AB}+\vv{AC}=\vv{BC}$}
		 {$\vv{AB}-\vv{BC}=\vv{AC}$}
	\loigiai{ 
		Ta có  $\vv{AB}+\vv{BC}+\vv{CA}=\vv{AC}+\vv{CA}=\vv{AA}=\vv{0}$.
	 }
\end{ex}

% nguồn: MaHoaNH_TracNghiem_Lop10
\begin{ex}%[0H5H2-2]%[ma:853c88cc]
	Cho ba điểm $A$, $B$, $C$ phân biệt.  Đẳng thức nào sau đây là đẳng thức \textbf{sai}?
	\choice 
		 {$\overrightarrow{BA}+\overrightarrow{AC}=\overrightarrow{BC}$}
		{$\overrightarrow{AB}+\overrightarrow{BC}=\overrightarrow{AC}$} 
		{\True $\overrightarrow{CA}+\overrightarrow{AB}=\overrightarrow{BC}$} 
		{$\overrightarrow{AB}-\overrightarrow{AC}=\overrightarrow{CB}$} 
	\loigiai{ 
	Theo quy tắc  ba điểm, ta có $\overrightarrow{CA}+\overrightarrow{AB}=\overrightarrow{CB}$ nên $\overrightarrow{CA}+\overrightarrow{AB}=\overrightarrow{BC}$ là sai.}
\end{ex}

% nguồn: MaHoaNH_TracNghiem_Lop10
\begin{ex}%[0H5H2-2]%[ma:a3c4c516]
	Đẳng thức nào sau đây  luôn đúng với mọi điểm $A$, $B$, $C$ bất kì?
	\choice 
		{$\overrightarrow{AB}+\overrightarrow{CB}=\overrightarrow{AC}$} 
		 {\True $\overrightarrow{CB}+\overrightarrow{AC}=\overrightarrow{AB}$}
		{$\overrightarrow{AB}-\overrightarrow{AC}=\overrightarrow{BC}$} 
		 {$\overrightarrow{BC}-\overrightarrow{AC}=\overrightarrow{AB}$}
	 \loigiai{
	Ta  có $\overrightarrow{CB}+\overrightarrow{AC}=\overrightarrow{AC}+\overrightarrow{CB}=\overrightarrow{AB}$.}
\end{ex}

% nguồn: MaHoaNH_TracNghiem_Lop10
\begin{ex}%[0H5H2-2]%[ma:71174e58]
	Chọn khẳng định  \textbf{sai}.
	\choice 
		 {Nếu $I$ là trung điểm đoạn $AB$ thì $\overrightarrow{AI}+\overrightarrow{IB}=\overrightarrow{AB}$}
		{\True Nếu $I$ là trung điểm đoạn $AB$  thì $\overrightarrow{IA}+\overrightarrow{BI}=\overrightarrow{0}$}
		{Nếu  $I$ là trung điểm đoạn $AB$ thì $\overrightarrow{AI}+\overrightarrow{BI}=\overrightarrow{0}$}
		{Nếu $I$ là trung  điểm đoạn $AB$ thì $\overrightarrow{IA}+\overrightarrow{IB}=\overrightarrow{0}$}
	 \loigiai{
		Ta có  $\overrightarrow{IA}+\overrightarrow{BI}=\overrightarrow{BI}+\overrightarrow{IA}=\overrightarrow{BA}$.\\
	Vậy $\overrightarrow{IA}+\overrightarrow{BI}=\overrightarrow{0}$  là sai.}
\end{ex}

% nguồn: MaHoaNH_TracNghiem_Lop10
\begin{ex}%[0H5H2-2]%[ma:44d45205]
	 Gọi $G$ là trọng tâm của tam giác $ABC$ và $M$ là điểm bất kỳ. Chọn khẳng định \textbf{sai} trong các khẳng định sau.
	\choice 
		{$\overrightarrow{AG}+\overrightarrow{BG}+\overrightarrow{CG}=\overrightarrow{0}$} 
		{$\overrightarrow{MA}+\overrightarrow{MB}+\overrightarrow{MC}=3\overrightarrow{MG}$} 
		 {$\overrightarrow{GA}+\overrightarrow{GB}+\overrightarrow{GC}=\overrightarrow{0}$}
		{\True  $\overrightarrow{MA}+\overrightarrow{MB}+\overrightarrow{MC}=\overrightarrow{MG}$}
	 \loigiai{
		Với $G$ là trọng tâm của tam  giác $ABC$ và $M$ là điểm bất kỳ, ta có
		\begin{itemize}
			 \item $\overrightarrow{AG}+\overrightarrow{BG}+\overrightarrow{CG}=\overrightarrow{0}\Leftrightarrow-\overrightarrow{GA}-\overrightarrow{GB}-\overrightarrow{GC}=\overrightarrow{0}\Leftrightarrow\overrightarrow{GA}+\overrightarrow{GB}+\overrightarrow{GC}=\overrightarrow{0}$.
			 \item $\overrightarrow{GA}+\overrightarrow{GB}+\overrightarrow{GC}=\overrightarrow{0}$.
			\item $\overrightarrow{MA}+\overrightarrow{MB}+\overrightarrow{MC}=3\overrightarrow{MG}$. 
		\end{itemize}
	} 
\end{ex}

% nguồn: MaHoaNH_TracNghiem_Lop10
\begin{ex}%[0H5H2-2]%[ma:134043d9]
	Gọi $O$ là tâm hình bình hành $ABCD$. Đẳng thức nào  sau đây \textbf{sai}?
	\immini{  \choice
	{$\overrightarrow{AB}-\overrightarrow{AD}=\overrightarrow{DB}$} 
	{$\overrightarrow{BC}-\overrightarrow{BA}=\overrightarrow{DC}-\overrightarrow{DA}$} 
	{$\overrightarrow{OA}-\overrightarrow{OB}=\overrightarrow{CD}$} 
	{\True $\overrightarrow{OB}-\overrightarrow{OC}=\overrightarrow{OD}-\overrightarrow{OA}$}}{\begin{tikzpicture}[line join=round,line cap=round,>=stealth,scale=.75,font=\footnotesize]
		\foreach \x/\y/\n  in {0/0/A,1/2/B,4/0/D} \coordinate (\n) at (\x,\y);
		\coordinate (C)  at ($(B)+(D)-(A)$);
		\coordinate (O)  at ($(B)!0.5!(D)$);
		\draw (A)--(B)--(C)--(D)--(A)--(C) 
			(B)--(D); 
		\foreach \t/\g  in {A/-135,C/45,B/135,D/-45,O/-90}{
			 \draw[fill=black] (\t) circle (1pt) node[shift={(\g:8pt)}]{$ \t $};
		} 
	\end{tikzpicture}
	 }
	 \loigiai{
		Ta  có
		$\overrightarrow{OB}-\overrightarrow{OC}=\overrightarrow{OD}-\overrightarrow{OA}\Leftrightarrow\overrightarrow{CB}=\overrightarrow{AD}$.\\ 
	Đáp  án này sai vì $\overrightarrow{CB}=\overrightarrow{DA}$ (do $ABCD$ là hình bình hành).}
\end{ex}

% nguồn: MaHoaNH_TracNghiem_Lop10
\begin{ex}%[0H5H2-2]%[ma:7d5c7704]
	Cho hình bình hành $ABCD$ tâm $O$. Đẳng thức nào sau đây  đúng?
	 \choice
		{\True $\overrightarrow{OA}+\overrightarrow{OC}=\overrightarrow{0}$} 
		{$\overrightarrow{OA}+\overrightarrow{CO}=\overrightarrow{0}$} 
		 {$\overrightarrow{AO}+\overrightarrow{OC}=\overrightarrow{0}$}
		 {$\overrightarrow{OA}-\overrightarrow{OC}=\overrightarrow{0}$}
	\loigiai{ 
	 \immini{
		Vì $ABCD$ là hình bình hành nên $O$  là trung điểm của $AC$.\\
		Vậy  $\overrightarrow{OA}+\overrightarrow{OC}=\overrightarrow{0}$.
		}{\begin{tikzpicture}[line join=round,line cap=round,>=stealth,scale=.75,font=\footnotesize]
		\foreach \x/\y/\n  in {0/0/A,1/2/B,4/0/D} \coordinate (\n) at (\x,\y);
		\coordinate  (C) at ($(B)+(D)-(A)$);
		\coordinate  (O) at ($(B)!0.5!(D)$);
		\draw  (A)--(B)--(C)--(D)--(A)--(C)
			 (B)--(D);
		\foreach \t/\g  in {A/-135,C/45,B/135,D/-45,O/-90}{
			\draw[fill=black] (\t)  circle (1pt) node[shift={(\g:8pt)}]{$ \t $};
		} 
	\end{tikzpicture}
		 }
	 }
\end{ex}

% nguồn: MaHoaNH_TracNghiem_Lop10
\begin{ex}%[0H5H2-2]%[ma:51694929]
	Cho hình bình hành $ABCD$ có tâm  $O$. Khẳng định nào sau đây \textbf{sai}?
	 \choice
		{\True $\overrightarrow{OB}+\overrightarrow{OD}=\overrightarrow{BD}$} 
		 {$\overrightarrow{OA}+\overrightarrow{OC}=\overrightarrow{0}$}
		{$\overrightarrow{AB}=\overrightarrow{DC}$} 
		{$\overrightarrow{AB}+\overrightarrow{AD}=\overrightarrow{AC}$} 
	 \loigiai{
	\immini{ 
		Ta có $O$ là trung điểm của  $BD\Rightarrow\overrightarrow{OB}+\overrightarrow{OD}=\overrightarrow{0}$.\\
		Do đó khẳng định $\overrightarrow{OB}+\overrightarrow{OD}=\overrightarrow{BD}$ là  sai.
		}{\begin{tikzpicture}[line join=round,line cap=round,>=stealth,scale=.75,font=\footnotesize]
		\foreach \x/\y/\n in {0/0/A,1/2/B,4/0/D} \coordinate (\n) at (\x,\y); 
		\coordinate  (C) at ($(B)+(D)-(A)$);
		\coordinate (O) at ($(B)!0.5!(D)$); 
		 \draw (A)--(B)--(C)--(D)--(A)--(C)
			 (B)--(D);
		\foreach \t/\g in  {A/-135,C/45,B/135,D/-45,O/-90}{
			\draw[fill=black] (\t) circle (1pt) node[shift={(\g:8pt)}]{$ \t  $};
		} 
	\end{tikzpicture}
		} 
	} 
\end{ex}

% nguồn: MaHoaNH_TracNghiem_Lop10
\begin{ex}%[0H5H2-2]%[ma:d9291e22]
	Cho lục giác đều $ABCDEF$ và $O$  là tâm của nó. Đẳng thức nào dưới đây là đẳng thức \textbf{sai}?
	 \choice
		 {$\overrightarrow{OA}+\overrightarrow{OC}+\overrightarrow{OE}=\overrightarrow{0}$}
		 {$\overrightarrow{BC}+\overrightarrow{FE}=\overrightarrow{AD}$}
		 {$\overrightarrow{OA}+\overrightarrow{OB}+\overrightarrow{OC}=\overrightarrow{EB}$}
		{\True $\overrightarrow{AB}+\overrightarrow{CD}+\overrightarrow{FE}=\overrightarrow{0}$} 
	\loigiai{ 
	\immini{ 
		 Ta có $\overrightarrow{AB}+\overrightarrow{CD}+\overrightarrow{FE}=\overrightarrow{AB}+\overrightarrow{BO}+\overrightarrow{AO}=2\overrightarrow{AO}\ne\overrightarrow{0}$.\\
		Do đó mệnh  đề $\overrightarrow{AB}+\overrightarrow{CD}+\overrightarrow{FE}=\overrightarrow{0}$ là sai.
		}{\begin{tikzpicture}[line join=round,line cap=round,>=stealth,scale=.75,font=\footnotesize]
		 \coordinate (O) at (0,0);
		\foreach \n/\p in {A/120,B/60,C/0,D/-60,E/-120,F/-180} \coordinate (\n) at ($(O)+({2*cos (\p)},{2*sin (\p)})$); 
		\draw  (A)--(B)--(C)--(D)--(E)--(F)--(A)--(D)
			(B)--(E) 
			 (C)--(F);
		\foreach \t/\g in {A/120,B/60,C/0,D/-60,E/-120,F/-180,O/-90}{ 
			\draw[fill=black] (\t) circle (1pt) node[shift={(\g:8pt)}]{$ \t  $};
		 }
	\end{tikzpicture}
		 }
	} 
\end{ex}

% nguồn: MaHoaNH_TracNghiem_Lop10
\begin{ex}%[0H5H2-2]%[ma:cc7f077b]
	Gọi $O$ là tâm hình bình hành $ABCD$; hai điểm $E$, $F$ lần  lượt là trung điểm $AB$, $BC$. Đẳng thức nào sau đây \textbf{sai}?
	\choice 
		{$\overrightarrow{DO}=\overrightarrow{EB}-\overrightarrow{EO}$} 
		 {$\overrightarrow{OC}=\overrightarrow{EB}+\overrightarrow{EO}$}
		{$\overrightarrow{OA}+\overrightarrow{OC}+\overrightarrow{OD}+\overrightarrow{OE}+\overrightarrow{OF}=\vec{0}$} 
		 {\True $\overrightarrow{BE}+\overrightarrow{BF}-\overrightarrow{DO}=\vec{0}$}
	\loigiai{ 
	\immini{ 
		Ta có $OF$, $OE$ lần lượt là đường trung bình của tam  giác $BCD$ và $ABC$.\\
		Suy ra $BEOF$ là hình bình  hành.\\
		 Khi đó $\overrightarrow{BE}+\overrightarrow{BF}=\overrightarrow{BO}$.\\
		Suy ra $\overrightarrow{BE}+\overrightarrow{BF}-\overrightarrow{DO}=\overrightarrow{BO}-\overrightarrow{DO}=\overrightarrow{OD}-\overrightarrow{OB}=\overrightarrow{BD}$. 
		}{\begin{tikzpicture}[line join=round,line cap=round,>=stealth,scale=1,font=\footnotesize]
		\foreach \x/\y/\n in {0/0/A,3/0/D,-1/-2/B} \coordinate (\n)  at (\x,\y);
		\coordinate (C) at ($(B)+(D)-(A)$); 
		 \coordinate (E) at ($(A)!1/2!(B)$);
		 \coordinate (F) at ($(B)!1/2!(C)$);
		\coordinate (O) at ($(A)!1/2!(C)$); 
		\draw  (A)--(B)--(C)--(D)--(A)--(C)
			(B)--(D) 
			(O)--(E)--(F)--(O); 
		 \foreach \t/\g in {A/90,B/-135,C/-45,D/45,E/180,F/-90,O/90}{
			 \draw[fill=black] (\t) circle (1pt) node[shift={(\g:8pt)}]{$ \t $};
		} 
	\end{tikzpicture}
		} 
	 }
\end{ex}

% nguồn: MaHoaNH_TracNghiem_Lop10
\begin{ex}%[0H5H2-2]%[ma:95d928f3]
	Cho hình vuông $ABCD$, có $M$  là giao điểm của hai đường chéo. Trong các mệnh đề sau, tìm mệnh đề \textbf{sai}?
	\choice 
		 {$\overrightarrow{AB}+\overrightarrow{BC}=\overrightarrow{AC}$}
		 {$\overrightarrow{AB}+\overrightarrow{AD}=\overrightarrow{AC}$}
		{$\overrightarrow{BA}+\overrightarrow{BC}=2\overrightarrow{BM}$} 
		 {\True $\overrightarrow{MA}+\overrightarrow{MB}=\overrightarrow{MC}+\overrightarrow{MD}$}
	 \loigiai{
	 \immini{
		Ta  có
		$$\begin{aligned}
		 &\overrightarrow{MA}+\overrightarrow{MB}=\overrightarrow{MC}+\overrightarrow{MD}\\
		\Leftrightarrow\ & \overrightarrow{MA}-\overrightarrow{MC}=\overrightarrow{MD}-\overrightarrow{MB}\\ 
		 \Leftrightarrow\ & \overrightarrow{CA}=\overrightarrow{BD}.
	\end{aligned}$$
	Vì  $\overrightarrow{CA}=\overrightarrow{BD}$ sai nên $\overrightarrow{MA}+\overrightarrow{MB}=\overrightarrow{MC}+\overrightarrow{MD}$ là sai.
	}{\begin{tikzpicture}[line join=round,line cap=round,>=stealth,scale=.75,font=\footnotesize]
		\foreach \x/\y/\n in {0/0/A,0/3/B,3/0/D} \coordinate (\n) at  (\x,\y);
		\coordinate (C) at ($(B)+(D)-(A)$); 
		\coordinate (M)  at ($(B)!0.5!(D)$);
		 \draw (A)--(B)--(C)--(D)--(A)--(C)
			(B)--(D); 
		\foreach  \t/\g in {A/-135,C/45,B/135,D/-45,M/-90}{
			\draw[fill=black] (\t)  circle (1pt) node[shift={(\g:8pt)}]{$ \t $};
		} 
	\end{tikzpicture}
	 }
 }
\end{ex}

% nguồn: MaHoaNH_TracNghiem_Lop10
\begin{ex}%[0H5H2-2]%[ma:4bbbd711]
	 Cho hình bình hành $ABCD$. Gọi $G$ là trọng tâm của tam giác $ABC$. Mệnh đề nào sau đây đúng?
	 \choice
		 {\True $\overrightarrow{GA}+\overrightarrow{GC}+\overrightarrow{GD}=\overrightarrow{BD}$}
		 {$\overrightarrow{GA}+\overrightarrow{GC}+\overrightarrow{GD}=\overrightarrow{CD}$}
		{$\overrightarrow{GA}+\overrightarrow{GC}+\overrightarrow{GD}=\overrightarrow{O}$} 
		{$\overrightarrow{GA}+\overrightarrow{GD}+\overrightarrow{GC}=\overrightarrow{CD}$} 
	\loigiai{ 
	 \immini{
		Vì $G$ là trọng tâm  của $\triangle ABC$ nên $\overrightarrow{GA}+\overrightarrow{GB}+\overrightarrow{GC}=\overrightarrow{0}$.\\
		 Do đó
		$$\begin{aligned}
		\overrightarrow{GA}+\overrightarrow{GC}+\overrightarrow{GD}&=\overrightarrow{GA}+\overrightarrow{GC}+\left(\overrightarrow{GB}+\overrightarrow{BC}+\overrightarrow{CD}\right)\\ 
		 &=\left(\overrightarrow{GA}+\overrightarrow{GB}+\overrightarrow{GC}\right)+\overrightarrow{BC}+\overrightarrow{CD}\\
		 &=\overrightarrow{BC}+\overrightarrow{CD}=\overrightarrow{BD}.
	\end{aligned}$$
	}{\begin{tikzpicture}[line join=round,line cap=round,>=stealth,scale=1,font=\footnotesize]
		\foreach \x/\y/\n in  {0/0/A,3/0/D,-1/-2/B} \coordinate (\n) at (\x,\y);
		\coordinate (C) at ($(B)+(D)-(A)$); 
		\coordinate (M) at ($(B)!1/2!(C)$); 
		\coordinate (O) at  ($(A)!1/2!(C)$);
		\coordinate (G)  at ($(B)!2/3!(O)$);
		 \draw (A)--(B)--(C)--(D)--(A)--(C)
			 (B)--(D)
			 (A)--(M);
		\foreach  \t/\g in {A/90,B/-135,C/-45,D/45,G/-30,M/-90,O/90}{
			\draw[fill=black] (\t) circle (1pt) node[shift={(\g:8pt)}]{$ \t $}; 
		} 
	\end{tikzpicture}
	} 
} 
\end{ex}


%%%%%%%%%%%%%%%% ĐÚNG/SAI %%%%%%%%%%%%%%%%

% nguồn: MaHoaNH_DungSai_Lop10
\begin{ex}%[0H5N2-2]%[ma:888353b1]
	Cho hình bình hành $ABCD$  tâm $O$. Hãy tìm các khẳng định \textbf{đúng}.
	 \choiceTF
		{\True $\vec{AB}+\vec{AD}=\vec{AC}$} 
		{$\vec{BA}+\vec{BD}=\vec{BC}$} 
		{\True $\vec{DA}=\vec{CB}$} 
		{\True $\vec{OA}+\vec{OB}+\vec{OC}\vec{OD}=\vec{0}$} 
	\loigiai{ 
		\begin{itemchoice}
			\itemch Vì $ABCD$  là hình bình hành nên$\vec{AB}+\vec{AD}=\vec{AC}$.
			\itemch $\vec{BA}+\vec{BD}=\vec{BC}$ là khẳng  định sai.
			\itemch  Vì $ABCD$ là hình bình hành nên $\vec{DA}=\vec{CB}$.
			\itemch Vì $ABCD$ là hình bình  hành nên $\vec{OA}+\vec{OB}+\vec{OC}\vec{OD}=\vec{0}$.
		\end{itemchoice}
	 }
\end{ex}

% nguồn: MaHoaNH_DungSai_Lop10
\begin{ex}%[0H5H2-2]%[ma:dd7c71f9]
	Gọi $M$ là trung điểm của đoạn thẳng $AB$.  Hãy tìm các khẳng định \textbf{đúng}.
	\choiceTF 
		 {$\vec{MA}=\vec{MB}$}
		{\True $\vec{AB}$ và $\vec{AM}$ cùng  phương}
		{$\vec{AB}$ và $\vec{MB}$ ngược  hướng}
		{\True $\left|\vec{AM}\right|=\left|\vec{BM}\right|$} 
	\loigiai{ 
		\begin{itemize}
			\item Do $M$ là trung điểm của đoạn thẳng $AB$ nên $\vec{AB}$ và  $\vec{AM}$ cùng phương.
			\item Do $M$ là trung điểm của đoạn thẳng $AB$  nên $\left|\vec{AM}\right|=\left|\vec{BM}\right|$.
		\end{itemize}
	 }
\end{ex}

% nguồn: MaHoaNH_DungSai_Lop10
\begin{ex}%[0H5H2-2]%[ma:af86b9e2]
	Cho hình bình hành $ABCD$. Tìm các  đẳng thức đúng trong các đẳng thức sau.
	\choiceTF 
		{\True $\vec{BA}+\vec{BC}=\vec{BD}$} 
		{\True $\vec{AD}=\vec{BC}$} 
		 {$\vec{AB}=\vec{CD}$}
		{\True  $\vec{AB}+\vec{AD}=\vec{AC}$}
	 \loigiai{
		\begin{itemchoice}
			\itemch $\vec{BA}+\vec{BC}=\vec{BD}$. 
			\itemch $ABCD$ là hình bình  hành nên $\vec{AD}=\vec{BC}$.
			\itemch $\vec{AB}=\vec{CD}$  thì $ABDC$ là hình bình hành.
			 \itemch $\vec{AB}+\vec{AD}=\vec{AC}$ (quy tắc hình bình hành).
		\end{itemchoice}
	 }
\end{ex}

% nguồn: MaHoaNH_DungSai_Lop10
\begin{ex}%[0H5H2-2]%[ma:01eb57bb]
	Cho ba điểm bất kỳ $A$, $B$, $C$. Hãy tìm các khẳng  định \textbf{đúng}.
	 \choiceTF
		{\True $\vec{CA} = \vec{BA}-\vec{BC}$} 
		{$\vec{AB} = \vec{AC}-\vec{CB}$} 
		 {\True $\vec{BC} = \vec{AC}+\vec{BA}$}
		 {\True $\vec{AB} + \vec{BC}=- \vec{CA}$}
	 \loigiai{
		\begin{itemchoice}
			\itemch  $\vec{CA} = \vec{BA}-\vec{BC}$.
			\itemch $\vec{AB} - \vec{AC}=\vec{CB} \Rightarrow  \vec{AB} = \vec{AC}+\vec{CB}$.
			\itemch $\vec{BC} -  \vec{BA}=\vec{AC} \Rightarrow \vec{BC} = \vec{AC}+\vec{BA}$.
			\itemch $\vec{AB} + \vec{BC}=\vec{AC}=- \vec{CA}$. 
		\end{itemchoice}
	 }
\end{ex}

% nguồn: MaHoaNH_DungSai_Lop10
\begin{ex}%[0H5H2-2]%[ma:3d324ad2]
	 Hãy cho biết tính đúng sai của mỗi đẳng thức sau.
	\choiceTF 
		{\True  $\overrightarrow{AB}+\overrightarrow{BC}=\overrightarrow{AC}$}
		{$\overrightarrow{AB}-\overrightarrow{AC}=\overrightarrow{BC}$} 
		{\True $\overrightarrow{AB}+\overrightarrow{BA}=\overrightarrow{0}$} 
		{$\overrightarrow{AB}-\overrightarrow{BC}=\overrightarrow{AC}$} 
	 \loigiai{
	\begin{itemchoice}
		\itemch  $\overrightarrow{AB}+\overrightarrow{BC}=\overrightarrow{AC}$.
		\itemch $\overrightarrow{AB}-\overrightarrow{AC}=\overrightarrow{CB}$. 
		\itemch $\overrightarrow{AB}+\overrightarrow{BA}=\overrightarrow{0}$ 
		\itemch  $\overrightarrow{AB}-\overrightarrow{BC}\ne \overrightarrow{AC}$.
	\end{itemchoice}
	 }
\end{ex}

% nguồn: MaHoaNH_DungSai_Lop10
\begin{ex}%[0H5H2-2]%[ma:c1859d95]
	Cho tam giác $ABC$. Các điểm $M$,  $N$, $P$ lần lượt là trung điểm của $AB$, $AC$, $BC$.
	 \choiceTF
		 {\True $\overrightarrow{A M}-\overrightarrow{A N}=\overrightarrow{NM}$}
		 {$\overrightarrow{MN}-\overrightarrow{NC}=\overrightarrow{MP}$}
		{\True $\overrightarrow{MN}-\overrightarrow{PN}=\overrightarrow{MP}$} 
		{$\overrightarrow{BP}-\overrightarrow{CP}=\overrightarrow{PC}$} 
	\loigiai{ 
		Dễ  thấy tứ giác $AMPN$, $MNCP$ là hình bình hành.
		\begin{center}
			\begin{tikzpicture}[line join=round,line cap=round,>=stealth,scale=1,font=\footnotesize]
					\foreach \x/\y/\n in {0/0/B,4/0/C,1/2/A}  \coordinate (\n) at (\x,\y);
					\coordinate (M) at ($(A)!1/2!(B)$); 
					 \coordinate (N) at ($(A)!1/2!(C)$);
					\coordinate (P) at  ($(B)!1/2!(C)$);
					\draw (A)--(B)--(C)--(A) 
						 (M)--(N)--(P)--(M);
					\foreach  \t/\g in {A/90,B/-135,C/-45,M/180,N/30,P/-90}{
						\draw[fill=black] (\t) circle (1pt) node[shift={(\g:8pt)}]{$ \t $}; 
					} 
				\end{tikzpicture}
		\end{center}
		\begin{itemchoice}
			 \itemch Vì $\overrightarrow{AM}-\overrightarrow{A N}=\overrightarrow{NM}$.
			 \itemch Vì $\overrightarrow{MN}-\overrightarrow{NC}=\overrightarrow{MN}-\overrightarrow{MP}=\overrightarrow{PN}\ (\overrightarrow{NC}=\overrightarrow{MP})$.
			\itemch Vì  $\overrightarrow{MN}-\overrightarrow{PN}=\overrightarrow{MN}+\overrightarrow{NP}=\overrightarrow{MP}$.
			 \itemch Vì $\overrightarrow{BP}-\overrightarrow{CP}=\overrightarrow{BP}+\overrightarrow{PC}=\overrightarrow{BC}$.
		\end{itemchoice}
	} 
\end{ex}

% nguồn: MaHoaNH_DungSai_Lop10
\begin{ex}%[0H5V2-2]%[ma:b96fcd46]
	Cho hình bình  hành $ABCD$ tâm $O$. Khi đó
	\choiceTF 
		{\True $\overrightarrow{BA}+\overrightarrow{AB}=\vec{0}$} 
		{\True $\overrightarrow{OA}+\overrightarrow{AC}=\overrightarrow{OC}$} 
		 {$\overrightarrow{OA}-\overrightarrow{OB}+\overrightarrow{DC}=\overrightarrow{CB}$}
		{\True $\overrightarrow{OB}+\overrightarrow{OA}-\overrightarrow{CA}=\overrightarrow{DC}$} 
	 \loigiai{
		\begin{itemchoice}
			\itemch Vì  $\overrightarrow{BA}+\overrightarrow{AB}=\overrightarrow{BB}=\vec{0}$.
			\itemch Vì $\overrightarrow{OA}+\overrightarrow{AC}=\overrightarrow{OC}$ (quy tắc ba  điểm).
			 \itemch Vì $\overrightarrow{OA}-\overrightarrow{OB}+\overrightarrow{DC}=\overrightarrow{OA}+\overrightarrow{BO}+\overrightarrow{DC}=\overrightarrow{BO}+\overrightarrow{OA}+\overrightarrow{AB}=\overrightarrow{BA}+\overrightarrow{AB}=\overrightarrow{0}$.
			\itemch Vì  $\overrightarrow{OB}+\overrightarrow{OA}-\overrightarrow{CA}=\overrightarrow{OB}+\overrightarrow{OA}+\overrightarrow{AC}=\overrightarrow{OB}+\overrightarrow{OC}=\overrightarrow{DO}+\overrightarrow{OC}=\overrightarrow{DC}$.
		\end{itemchoice}
	 }
\end{ex}


%%%%%%%%%%%%%%%% TỰ LUẬN %%%%%%%%%%%%%%%%

% nguồn: MaHoaNH_TuLuan_Lop10
\begin{bt}%[0H5H2-2]%[ma:2cebee43]
	Cho bốn điểm bất  kỳ $A$, $B$, $C$ và $D$. Hãy chứng minh đẳng thức $\overrightarrow{AB}+\overrightarrow{CD}=\overrightarrow{AD}+\overrightarrow{CB}$.
	\loigiai{ 
		\begin{eqnarray*}
			 \overrightarrow{AB}+\overrightarrow{CD}&=&\left(\overrightarrow{AD}+\overrightarrow{DB}\right)+\left(\overrightarrow{CB}+\overrightarrow{BD}\right)\\
			 &=&\overrightarrow{AD}+\overrightarrow{CB}+\left(\overrightarrow{BD}+\overrightarrow{DB}\right)\\
			 &=&\overrightarrow{AD}+\overrightarrow{CB}+\overrightarrow0\\
			&=&\overrightarrow{AD}+\overrightarrow{CB}. 
		\end{eqnarray*}
	 }
\end{bt}

% nguồn: MaHoaNH_TuLuan_Lop10
\begin{bt}%[0H5V2-2]%[ma:b2c36403]
	Cho sáu điểm $A$, $B$, $C$, $D$, $E$, $F$. Chứng minh rằng 
	\begin{enumerate}
		\item  $\overrightarrow{AB}+\overrightarrow{CD}+\overrightarrow{EF}-\overrightarrow{CB}-\overrightarrow{ED}=\overrightarrow{AF}$.
		 \item $\overrightarrow{AB}-\overrightarrow{AF}+\overrightarrow{CD}-\overrightarrow{CB}+\overrightarrow{EF}=\overrightarrow{ED}$.
		\item $\overrightarrow{AB}+\overrightarrow{CD}=\overrightarrow{AD}+\overrightarrow{CB}$. 
		 \item $\overrightarrow{AC}+\overrightarrow{BD}+\overrightarrow{EF}=\overrightarrow{A F}+\overrightarrow{BC}+\overrightarrow{ED}$.
	\end{enumerate}
	\loigiai{ 
		\begin{enumerate}
			 \item
			$$\begin{aligned}
			 \overrightarrow{AB}+\overrightarrow{CD}+\overrightarrow{EF}-\overrightarrow{CB}-\overrightarrow{ED}
			&=\overrightarrow{AB}+\overrightarrow{CD}+\overrightarrow{EF}+\overrightarrow{BC}+\overrightarrow{DE}\\ 
			&=\left(\overrightarrow{AB}+\overrightarrow{BC}\right)+\left(\overrightarrow{CD}+\overrightarrow{DE}\right)+\overrightarrow{EF}\\ 
			&=\overrightarrow{AC}+\overrightarrow{CE}+\overrightarrow{EF}=\overrightarrow{AE}+\overrightarrow{EF}=\overrightarrow{AF}. 
		\end{aligned}$$
		\item  $\overrightarrow{AB}-\overrightarrow{AF}+\overrightarrow{CD}-\overrightarrow{CB}+\overrightarrow{EF}=\overrightarrow{FB}+\overrightarrow{BD}+\overrightarrow{EF}=\overrightarrow{FD}+\overrightarrow{EF}=\overrightarrow{ED}$.
		 \item $\overrightarrow{AB}+\overrightarrow{CD}=\overrightarrow{AD}+\overrightarrow{CB}\Leftrightarrow\overrightarrow{AB}-\overrightarrow{AD}=\overrightarrow{CB}-\overrightarrow{CD}\Leftrightarrow\overrightarrow{DB}=\overrightarrow{DB}$.
		\item 
		$$\begin{aligned}
		&\overrightarrow{AC}+\overrightarrow{BD}+\overrightarrow{EF}=\overrightarrow{AF}+\overrightarrow{BC}+\overrightarrow{ED}\\ 
		\Leftrightarrow\ &  \left(\overrightarrow{AC}-\overrightarrow{AF}\right)+\left(\overrightarrow{BD}-\overrightarrow{BC}\right)+\left(\overrightarrow{EF}-\overrightarrow{ED}\right)=\vec{0}\\
		 \Leftrightarrow\ &\overrightarrow{FC}+\overrightarrow{CD}+\overrightarrow{DF}=\vec{0}\Leftrightarrow\overrightarrow{FD}+\overrightarrow{DF}=\vec{0}.
	\end{aligned}$$
\end{enumerate}}
\end{bt}
